\section{From a literature graph to a testable banking hypothesis}
\label{exec20:hypothesis-scope}
A bank translating a business question into a causal analysis needs more than a plausible arrow. Suppose a servicing team asks whether a digital reminder reduces missed payments among customers whose income has become irregular. Literature can suggest mechanisms involving attention, financial stress, access and payment timing. A useful assistant should expose those possibilities and their sources, then help specify an intervention, eligible population, outcome horizon and decision clock. It should not quietly turn a sentence about association into an identified effect of sending a reminder. Tong and colleagues provide a concrete design for using language models and a literature graph to generate research hypotheses. Their psychological setting is informative about the design of that assistance, while supplying no validation of the banking intervention.~\citep{exec20_hypothesis_generation}

\subsection{What the literature graph represents}
\label{exec20:hypothesis-claims}
The inspected July 2024 article describes three stages: retrieve literature, extract relation pairs and generate hypotheses from predicted links. From a wider psychological and neuroscience collection, the authors select 43,312 articles using a journal-title criterion; 35,797 come from \emph{Frontiers in Psychology}. That concentration is a property of the selected evidence, not a census of psychological knowledge. PDF cleaning and chunks of at most 4,000 tokens precede GPT-4 extraction. The prompt explicitly requests both causal and correlational relations, their concepts, and the sign of causal relations. Model filtering and periodic expert feedback follow. The resulting graph is reported to contain about 197,000 concepts and 235,000 connections. Neo4j supplies storage; node2vec embeddings and neighbor Jaccard similarity guide candidate-link selection. GPT-4 then turns selected concept pairs into hypotheses. These are the authors' reported components and counts, not a reproduced implementation.~\citep{exec20_hypothesis_generation}

A bank-facing evidence object should retain more information than the concept pair. Let $q$ identify a source claim, $d_q$ its document and version, $s_q$ the inspected passage, $r_q$ its asserted relation type, and $c_q$ its population, timing and study-design context. A useful proposed representation is
\begin{equation}
  e_q=(d_q,s_q,r_q,c_q,\hbox{extraction status},\hbox{review status}).
  \label{exec20:hypothesis-evidence}
\end{equation}
The claim graph constructed from these objects need not be an acyclic structural causal model. Different studies can use different populations, variable definitions and time resolutions; an article can report a causal interpretation whose identifying assumptions remain disputed. Deduplication must preserve disagreement and source independence. A collection of paraphrases of one study cannot acquire the status of several replications. Model confidence and an expert's agreement that the text contains a claim are also different quantities from evidence that the claim is true in the proposed bank population.

The paper's extraction check makes this distinction measurable. Among 238 pairs classified as causal by GPT-4, the main article reports 156 agreed causal labels and 31 pairs not recognized as relationships. The complementary 51 therefore have a different recognized relation label under that classification. Consequently,
\begin{equation}
 \frac{156}{238}\simeq 65.55\%,\qquad
 \frac{207}{238}\simeq 86.97\%,\qquad
 \frac{31}{238}\simeq 13.03\%.
 \label{exec20:hypothesis-denominators}
\end{equation}
The small rounding differences from the article's displayed percentages do not change the substantive point: relation-existence agreement is a different endpoint from causal-label agreement. Neither denominator measures an intervention's ground-truth effect. The authors' familiar-paper evaluation and feedback process supplies useful exploratory evidence; it does not establish performance on previously unseen banking mechanisms.~\citep{exec20_hypothesis_generation,exec20_hypothesis_reference}

\subsection{Why a predicted link remains a hypothesis}
\label{exec20:hypothesis-graph}
For neighbor sets $N(a)$ and $N(b)$, the usual Jaccard score is
\begin{equation}
 J(a,b)=\frac{|N(a)\cap N(b)|}{|N(a)\cup N(b)|}.
 \label{exec20:hypothesis-jaccard}
\end{equation}
This expresses shared graph context under a chosen neighbor convention. A denominator of zero needs an explicit implementation policy. Direction, edge types, duplicate claims and temporal aggregation also affect what a neighborhood means. A score in $[0,1]$ is not thereby a probability of causation, a confidence level or an effect magnitude.

An authored counterexample illustrates the missing implication. Take mutually independent zero-mean $Z,\varepsilon_A,\varepsilon_B,\varepsilon_P$ and the structural equations
\begin{equation}
 A=Z+\varepsilon_A,\qquad B=Z+\varepsilon_B,
 \qquad P=A+B+\varepsilon_P.
 \label{exec20:hypothesis-world}
\end{equation}
The DAG has $Z\to A$, $Z\to B$, $A\to P$ and $B\to P$. Using undirected adjacency neighborhoods, $N(A)=N(B)=\{Z,P\}$, so $J(A,B)=1$. Yet intervening on $A$ changes neither $Z$ nor the equation for $B$: the derivative of $E[B\mid\operatorname{do}(A=a)]$ with respect to $a$ is zero. Shared cause and shared consequence create identical neighborhoods without a direct or total effect of $A$ on $B$. This is a statement about the displayed construction; the uninspected package's neighbor convention is not assumed.~\citep{exec20_hypothesis_reference}

For a banking interpretation, let $A$ encode a reminder policy, $B$ a repayment response, and $Z$ a common predisposition in a hypothetical observational population. If a planner wrongly treats similarity one as a success probability for an intervention worth USD 100 at cost USD 20, the apparent margin is USD 80. Under the stipulated zero-effect law, the incremental margin is USD -20. No bank profit or package failure is measured here. The construction shows why a candidate generator must return a proposed relation with its evidence, while a subsequent identification step supplies any admissible causal effect or refusal.

\subsection{What the reported evaluation can establish}
\label{exec20:hypothesis-evaluation}
The authors focus on well-being and generate 130 hypotheses. The principal comparison uses 30 randomly selected model-and-graph hypotheses, 30 expert-selected hypotheses, 30 hypotheses sampled from 41 ideas supplied by 13 responding doctoral students, and 30 Claude-2 hypotheses selected by GPT-4 from 50 generated ideas. Three experts rate novelty and usefulness, with hypotheses presented in randomized order. These pipelines differ in generation and selection. The comparison therefore concerns the resulting pipelines and selected ideas, rather than isolating one graph component under identical resources. Reported reviewer correlations vary, including a novelty correlation of 0.069 for one reviewer pair; an average rating cannot erase that disagreement.~\citep{exec20_hypothesis_generation}

The paper reports higher novelty for its randomly selected graph-assisted ideas than for the Claude comparison. The expert-selected arm is not uniformly significant: Table 5 gives mean-comparison $p=0.085$ and median-comparison $p=0.550$ against Claude. A later 60-versus-60 GPT-4 ablation reports a novelty advantage but a usefulness comparison of $p=0.1937$. Nonsignificance does not prove equal usefulness, nor does an insignificant human-versus-model comparison establish equivalence to experts without a justified equivalence margin and suitable design. These endpoints are judgments of ideas; the proposed mechanisms were not validated by the reported rating experiment. BERT representations and t-SNE plots describe semantic structure and require their own interpretation. They do not replace an intervention experiment.~\citep{exec20_hypothesis_generation}

Several printed numerical statements also require reconciliation before the proposal inherits their inferential force. The usefulness paragraph pairs $F(3,116)=5.25$ with $p=0.553$. Independent evaluation of the ordinary $F$ upper tail gives approximately $0.001967$. In the semantic-distance analysis, the printed $F(3,1652)=84.1611$ is paired with explained variance 86.96\%. Under ordinary one-factor ANOVA with the printed degrees of freedom,
\begin{equation}
 R^2=\frac{d_1F}{d_1F+d_2}
     =\frac{3(84.1611)}{3(84.1611)+1652}
     \simeq 0.132573.
 \label{exec20:hypothesis-anova}
\end{equation}
These are arithmetic checks conditional on the stated conventional model, not a reconstruction of unavailable raw ratings. The correct repair is to obtain the analysis and resolve which printed quantities or definitions need correction. The prose also varies from its table for a human-comparison $t$ value and calls a median comparison significant where the table gives $p=0.550$. Hypotheses contributed by the same person, repeated ratings and distances sharing the same hypothesis create dependence questions; their inferential treatment cannot be recovered from a significance label alone.~\citep{exec20_hypothesis_generation,exec20_hypothesis_reference}

\subsection{A public evaluation before a bank prototype trial}
\label{exec20:hypothesis-prototype}
The proposed library can adopt the separation of literature extraction, graph-assisted candidate generation and expert evaluation as a design hypothesis. Direct code adoption remains premature: the selected article does not itself qualify immutable implementation identities, exact model snapshots, full prompts and thresholds, dependency licenses, or complete corpus redistribution rights. Its data statement makes only part of the evidence available with supplementary material and offers additional data on request. The earlier preprint's appendices are a different version and cannot silently substitute for the revised supplement. The article's own limitations acknowledge extraction error, synonymous concepts, a restricted hypothesis sample and a shift in scarce expert work from generation to evaluation.~\citep{exec20_hypothesis_generation}

A first public experiment should freeze a licensed corpus and evaluate four linked but separate questions. Can extraction recover source-supported claims, relation type, direction, sign and uncertainty on held-out documents? Can an assistant translate a business description into an operational population, intervention and outcome without inventing assignment or observability? Does graph-assisted generation improve blind ratings of useful, testable hypotheses relative to retrieval-only, language-model-only and human-assisted baselines under matched budgets? Finally, do proposals survive identification and intervention checks in generated banking mechanisms with known truth and in suitable publicly documented experiments? Novelty and semantic diversity remain secondary to correctness and decision loss when an action is proposed.

Document families, shared studies and generated mechanism families should be separated across development and evaluation to reduce leakage. The public protocol should retain contradicted claims, extraction failures, unavailable answers and refusal. Human evaluation needs a declared sampling unit, a predeclared treatment of reviewer disagreement and dependence, and a held-out adjudication process. Actual calls, retries, tokens, expert minutes and component compute should be recorded. The comparison should price the full review workflow rather than assume that automatic hypothesis production removes expert work.

Across library, CLI and MCP, the proposed output should distinguish \texttt{extracted\_claim}, \texttt{hypothesis}, \texttt{reviewed\_assumption} and \texttt{identified\_query}. It should carry immutable source anchors, event-time definitions, unresolved conflicts and the checks still required for promotion. A client requesting an effect must not receive a link score under an effect field. These contracts are a concrete route to a useful public prototype; qualification of a banking intervention still requires evidence about that intervention.

